\documentclass[11pt]{article}

% --- Core packages ---
\usepackage[utf8]{inputenc}
\usepackage[T1]{fontenc}
\usepackage{lmodern}
\usepackage[english]{babel}

% --- Math ---
\usepackage{amsmath}
\usepackage{amssymb}

% --- Tables, figures, layout ---
\usepackage{booktabs}
\usepackage{graphicx}
\usepackage{geometry}
\geometry{margin=1in}

% --- Citations (natbib with numeric style to match original look) ---
\usepackage[numbers,sort&compress]{natbib}

% --- Colors (required for blue!50!black mixing syntax in hyperref) ---
\usepackage{xcolor}

% --- Hyperlinks and URLs (load near last) ---
\usepackage{url}
\usepackage{hyperref}
\hypersetup{
  colorlinks=true,
  linkcolor=blue!50!black,
  citecolor=blue!50!black,
  urlcolor=blue!50!black,
}

% --- Convenience macros for cosmology notation ---
\newcommand{\LCDM}{$\Lambda$CDM}
\newcommand{\wCDM}{$w$CDM}
\newcommand{\wowaCDM}{$w_0 w_a$CDM}

% --- Document metadata ---
\title{The Avenridge Evaluation and Publication Framework:\\
Locked, Reproducible Evaluation with Calibration Assessment\\
and Null-Result Parity}
\author{Earl Dixon\thanks{The author is the sole operator of
Avenridge Institute, a research-methodology entity established to
develop and apply the framework described herein. The case studies
cited as Avenridge Institute technical reports are archived in the
public repository referenced in those citations and are not
peer-reviewed publications.}}
\date{\today}

\begin{document}
\maketitle

\begin{abstract}
We introduce the Avenridge Evaluation and Publication Framework
(AEPF), an operationalized methodology for evaluation-intensive
computational domains. AEPF integrates established
disciplines---clinical-trial pre-registration, high-energy-physics
blind analysis, machine-learning calibration assessment, and
open-science null reporting---into a single workflow with concrete
operational consequences: cryptographic anchoring of methodology to
a public commit chain, calibration assessment as a standing
pre-registration requirement, null-result publication parity,
constrained-interpretation reporting, and a pre-specification
viability gate that retires structurally over-determined studies. We
illustrate the framework with four case studies---spanning
cosmological residual analysis, transformer calibration auditing,
framing-robustness testing of an open-weight language model, and a
viability-gate disposition---together with a methods-level validation
of the calibration-audit instrument against synthetic ground truth.
The instrument validation establishes that the audit is sound: zero
false-acceptance of broken calibration across the tested
miscalibration families, so that a null verdict reported under AEPF
reflects the evaluated system rather than an artifact of the tool.
The novelty claimed is in the integration and artifact-level
operationalization, not in any individual component.
\end{abstract}

\section{Introduction}
\label{sec:introduction}

This paper introduces a methodology designed to constrain the
relationship between an evaluation question, the analytical
machinery applied to it, and the verdict that is ultimately
reported. The framework is anchored in cryptographic commit
discipline applied to pre-registered methodology, materialized
data sources, and published outcomes, so that a reader of an
AEPF-locked study can verify those relationships directly against
artifacts rather than infer them from the narrative of a paper.
In May 2026, the framework was applied to a cosmological audit
that returned a NULL verdict reported under two interpretive
frames \cite{dixon2026ersaf}; days later, it was applied as a
viability gate that retired a candidate computational study
before lock \cite{dixon2026dynametrixhrrr}. Both outcomes are
treated below as case studies of a single methodological
framework.

The cosmological audit just referenced is the most recent of a
sequence of applications. The framework has also been used across
atmospheric forecasting calibration \cite{dixon2026avenridge},
gravitational-wave residual analysis at GW150914
\cite{dixon2026spinphase}, and language-model calibration
assessment across multiple model architectures
\cite{dixon2026distilbert, dixon2026toxicbert,
dixon2026chatgptdetector}. The framework's central commitment is
that locking methodology before data observation, anchoring
artifacts cryptographically, and committing institutionally to
publishing null findings produces evaluations whose
results---positive or negative---are interpretable in ways that
benchmark-driven evaluation alone does not support.

A complementary use of the framework was demonstrated days after
the ERSAF result was published: as a structured \emph{viability
gate} preceding lock. A candidate study proposing to test whether
combining Dynametrix v3 calibrated severe-weather probabilities with
NWS High-Resolution Rapid Refresh (HRRR) forecasts would produce a
measurably better-calibrated joint prediction was subjected to a
pre-committed exploratory independence diagnostic before its
pre-registration was advanced to lock. The diagnostic measured the
Pearson correlation of the two predictors' residual errors at
$r = 0.9839$ (95\% bootstrap CI $[0.9816, 0.9860]$, $n = 2{,}016$
location-hours, 375 verified severe-weather events), far above the
pre-committed threshold for proceeding. The candidate pre-registration
was archived as ``considered but deferred'' rather than advanced to
lock \cite{dixon2026dynametrixhrrr}. The framework as articulated in
this paper handles full-lifecycle studies that lock, execute, and
report; the viability-gate pattern---a refinement that emerged in
operational practice and is described in
Section~\ref{subsec:viability-gate}---extends the lifecycle backward
to handle pre-specification dispositions.

This paper introduces the \textbf{Avenridge Evaluation and
Publication Framework (AEPF)}, an operationalized methodology for
evaluation-intensive computational domains. AEPF is not a new
theoretical proposal. It integrates established
disciplines---pre-registration from clinical trials and psychology
\cite{helsinki2013, nosek2018preregistration, chambers2013registered},
blind analysis protocols from high-energy physics \cite{klein2005blind},
calibration assessment from modern machine learning
\cite{brier1950, guo2017calibration, naeini2015ece}, and
reproducibility frameworks from the machine-learning community
\cite{pineau2021reproducibility}---into a single operationalizable
framework with specific artifact-level implementations. The novelty
AEPF claims is in the integration and operationalization, not in
any individual component.

\subsection{Problem statement}
\label{subsec:problem}

Modern applied computational systems increasingly operate in
high-consequence environments where predictive accuracy alone is
insufficient to establish trustworthiness. Systems may demonstrate
strong benchmark performance while remaining poorly calibrated,
insufficiently reproducible, vulnerable to evaluation contamination,
or susceptible to post-hoc interpretive bias.

These challenges are not unique to artificial intelligence.
Scientific measurement disciplines have historically relied on
methodological constraints intended to separate observation from
interpretation, preserve auditability, and reduce opportunities for
unconscious researcher influence. Yet evaluation methodologies
across applied AI and adjacent evaluation-intensive computational
domains frequently emphasize aggregate performance metrics while
treating calibration quality, artifact preservation, null-result
publication, and locked evaluation procedures as secondary
considerations.

This imbalance creates several specific risks:
\begin{itemize}
  \item Evaluation contamination through repeated exposure to fixed
    benchmark corpora, encouraging adaptation to benchmark
    characteristics rather than to underlying capability
    \cite{recht2019imagenet}.
  \item Over-interpretation of statistically fragile findings,
    particularly when multiple-testing corrections are applied
    inconsistently or post-hoc.
  \item Benchmark optimization without corresponding reliability
    improvements---high discriminative scores in the presence of
    poor calibration.
  \item Limited reproducibility arising from insufficient provenance
    tracking of data, code, computational environment, and
    methodology decisions.
  \item Selective publication favoring positive outcomes while
    under-reporting null findings, biasing the cumulative record
    toward apparent confirmation.
\end{itemize}

Calibration is particularly important because predictive systems
must not only distinguish correct from incorrect outcomes; they
must appropriately represent uncertainty. A system reporting
80\,\% confidence should empirically demonstrate correctness near
that rate over sufficient evaluation samples
\cite{guo2017calibration}. Systems exhibiting strong discrimination
while remaining poorly calibrated may encourage overconfidence and
reduce operational reliability---a failure mode documented in
modern neural network classifiers \cite{guo2017calibration,
minderer2021revisiting} and one we have observed empirically across
language-model calibration audits performed under the framework
introduced here \cite{dixon2026distilbert, dixon2026toxicbert,
dixon2026chatgptdetector}.

Reproducibility likewise depends on preserving evaluation artifacts,
recording methodological decisions, and constraining the surface
across which post-hoc modifications can occur. Without these
controls, evaluation procedures risk unintentionally rewarding
iterative adaptation to benchmark structure rather than measuring
genuine system capability.

\subsection{AEPF as an operationalized framework}
\label{subsec:aepf-overview}

AEPF is operationalized through specific artifacts and procedures,
not through abstract values. The framework prescribes:

\begin{enumerate}
  \item \textbf{Pre-registration documents} with locked methodology
    sections, committed before any analysis touches the evaluation
    data.
  \item \textbf{Cryptographic commit references} (SHA-256 lock
    commits with named version tags) anchoring the methodology to
    a specific reviewable state.
  \item \textbf{Materialization manifests} recording data-file
    SHA-256 hashes, software-environment version pinning, and
    per-step transformation provenance.
  \item \textbf{Calibration assessment} alongside discriminative
    metrics---typically Brier score \cite{brier1950}, Expected
    Calibration Error \cite{naeini2015ece}, and reliability diagrams
    reported as required disclosures.
  \item \textbf{Cross-model review} before lock commit, distinct
    from the operator's own review, producing numbered findings
    with severity classifications and explicit pass/fail
    dispositions.
  \item \textbf{Single-execution discipline}: the locked methodology
    runs once on observed data; re-runs are permitted only for
    reproducibility verification and do not produce new outcome
    documents.
  \item \textbf{Result documents with constrained interpretation
    sections} that, when multiple permissible interpretations apply
    to a result, document each frame side-by-side without claiming
    one stronger than the other given the locked methodology's
    scope.
  \item \textbf{Explicit publication of null findings} under the
    same lock-commit discipline as positive findings, with
    institutional commitment recorded prospectively in the
    pre-registration.
\end{enumerate}

\subsection{Scope and objective}
\label{subsec:scope}

AEPF does not prescribe scientific conclusions, model architectures,
optimization strategies, or domain-specific theories. The framework
is purely procedural: it specifies how evaluations are
pre-registered, locked, executed, and reported. Results obtained
under AEPF discipline make no stronger claims than the locked
methodology permits, and AEPF itself makes no claims about what
specific scientific or engineering decisions evaluators should
reach.

The objective is institutional evaluation discipline: improving the
interpretability, reproducibility, and auditability of computational
evaluation outcomes by making methodology decisions binding
artifacts rather than malleable practices. The work AEPF does for
computational evaluations is structurally analogous to what
Registered Reports \cite{chambers2013registered} do for empirical
psychology and what blind analysis protocols \cite{klein2005blind}
do for high-energy physics: it separates the design of an
evaluation from the observation of its outcome in a way that is
mechanically enforceable rather than purely aspirational.

\section{Related Challenges and Existing Frameworks}
\label{sec:related}

Several disciplines have established components of evaluation rigor
that AEPF integrates. This section briefly surveys each and
positions AEPF's contribution as integration and operationalization
rather than novel methodology.

\paragraph{Pre-registration.}
Prospective registration of methodology is well-developed in
clinical trials, where it is required by the Helsinki Declaration's
Article~35 \cite{helsinki2013} and operationalized through
ClinicalTrials.gov \cite{zarin2011clinicaltrials}. In psychology
and adjacent behavioral sciences, the Open Science Framework's
preregistration protocols \cite{nosek2018preregistration} and the
Registered Reports publication format \cite{chambers2013registered}
have extended the discipline to observational and quasi-experimental
studies. AEPF inherits the prospective-registration discipline from
this literature and adapts it to computational-evaluation contexts
where data, code, computational environment, and methodology must
be locked simultaneously rather than in separable stages.

\paragraph{Blind analysis.}
High-energy physics has applied blind analysis protocols for
decades, separating the determination of analysis cuts and event
selection from the observation of the signal region
\cite{klein2005blind}. The Large Hadron Collider experiments extend
this with pre-unblinding review processes that resemble AEPF's
cross-model review stage in function if not in specific form. AEPF
generalizes the blind-analysis principle---methodology locked
before signal observation---to computational evaluations where the
signal region is the evaluation corpus rather than a kinematic
phase space.

\paragraph{Calibration assessment.}
Probabilistic forecasting verification has a long history
\cite{brier1950}, but its systematic application to neural network
outputs is comparatively recent. Guo and coauthors documented
systematic overconfidence in modern neural classifiers
\cite{guo2017calibration} and proposed Expected Calibration Error
as a measurement standard, building on the binning approach of
Naeini and coauthors \cite{naeini2015ece}. Subsequent work has
explored calibration under distribution shift
\cite{minderer2021revisiting}. AEPF requires calibration assessment
as a standard reporting category, alongside whichever discriminative
metrics are appropriate to the domain.

\paragraph{Reproducibility frameworks.}
The machine-learning community has produced multiple
reproducibility-discipline artifacts in recent years, including the
NeurIPS reproducibility checklist \cite{pineau2021reproducibility}
and standardization efforts via Papers with Code. These frameworks
focus primarily on enabling reproduction of published results
through code-and-data release. AEPF extends the reproducibility
scope to include the methodology itself---pre-locked before
observation rather than reconstructed post-hoc from a paper
description.

\paragraph{Open science and null-result publication.}
The broader open-science movement \cite{nosek2015open} has
advocated for systematic publication of null results to address
publication bias \cite{ioannidis2005}. AEPF operationalizes this
commitment at the methodology-lock stage: the pre-registration
includes binding language that the result will be published
regardless of outcome direction, with the same artifact discipline
applied to NULL and POSITIVE outcomes.

\paragraph{The integration claim.}
AEPF's contribution is not any of these components in isolation,
each of which has substantial prior literature. The contribution is
a single operationalized framework that integrates prospective
registration, cryptographic anchoring, materialization manifests,
calibration assessment, cross-model review, single-execution
discipline, constrained-interpretation reporting, and null-result
publication into a methodology applicable across the heterogeneous
evaluation contexts that arise in applied artificial intelligence
and adjacent computational sciences. To our knowledge, no existing
framework combines all of these components in an artifact-level
operational form; AEPF is presented here as one such combination,
with case studies in subsequent sections demonstrating its
application across multiple domains.

\section{Framework Specification}
\label{sec:framework-spec}

This section describes the operational artifacts that compose AEPF in
sufficient detail to permit independent re-implementation. The
framework is organized around the temporal lifecycle of an evaluation:
pre-registration, lock-time commitment, execution, and reporting. Each
phase produces specific artifacts that constrain subsequent phases,
and each constraint is enforceable through reference to a
cryptographically dated commit rather than through institutional
trust.

\subsection{Pre-registration and cross-model review}
\label{subsec:prereg-review}

The pre-registration document is the central artifact of an AEPF
evaluation. It is finalized and committed before any analysis touches
the evaluation data. The document specifies the evaluation question
and the population to which it is addressed; the data sources, with
version pins or commit references for any external data; the baseline
models or comparison conditions, with their parameter values fixed
before unblinding; the observable quantities to be computed, including
preprocessing and transformation steps; the statistical procedures,
including significance thresholds, multiple-testing correction methods,
and any tolerance windows; the decision rules specifying the
conditions under which the result will be classified as POSITIVE,
NULL, or AMBIGUOUS; the required disclosures, including known
systematics and methodological limitations; and a publication
commitment binding the operator to publish the result regardless of
outcome direction.

Before lock, the pre-registration undergoes \emph{cross-model review}:
an independent review performed by a reviewer institutionally
separable from the operator. The reviewer produces a written review
containing numbered findings classified by severity (Critical,
Significant, Minor, Note) with explicit pass/fail dispositions per
finding. The operator is permitted to disagree with a reviewer finding
but must document the disagreement in writing and commit the
documentation alongside the pre-registration. A reviewer Critical
finding that the operator does not address blocks the lock commit.

Cross-model review is distinct from peer review at publication: it
operates on the pre-registered methodology, before data have been
observed, with the purpose of preventing avoidable methodological
errors from being locked into the artifact chain. Peer review at
publication remains valuable but cannot address pre-registration
errors after the methodology has been frozen.

\subsection{Materialization manifest}
\label{subsec:materialization}

External data sources are not loaded by URL, by package import, or by
any mechanism that could resolve to different bytes at different
times. Each external source is materialized: fetched once, written to
a project-local file, and recorded in a manifest file committed
alongside the pre-registration. The manifest specifies, per source:
the source identifier (URL, DOI, or commit reference of the upstream
repository); the local file path within the project; the SHA-256 hash
of the materialized file; the timestamp of materialization; and the
software environment used to perform the materialization.

Analysis code reads from the local materialized files, never from
upstream sources. If an upstream source updates between materialization
and execution, the manifest preserves the bytes that were actually
used; if the upstream is later retracted or modified, the locked
artifact chain remains independently verifiable. For the ERSAF
cosmological audit, the materialization manifest pinned the
Pantheon+ supernova compilation \cite{brout2022pantheon} to specific
file hashes, and the DESI DR1 BAO measurements
\cite{adame2024desi3, adame2024desi4} to file hashes derived from the
upstream \texttt{CobayaSampler/bao\_data} repository at a named
commit.

\subsection{Lock-time commitment}
\label{subsec:lock-commit}

When the pre-registration, cross-model review record, materialization
manifest, and analysis code are all complete, they are committed to
the version-control repository in a single lock commit. The lock
commit is then tagged with a signed git tag following the convention
\texttt{<framework>-v<version>-lock} (e.g., \texttt{ersaf-v1.0-lock}).
The tag is pushed to a public remote before any analysis touches the
evaluation data.

The lock commit establishes a single cryptographically dated reference
point. All subsequent commits to the repository are causally
downstream of the lock and are distinguishable from the pre-registered
state. The tag, once published to a public remote, cannot be moved
without leaving a record in the repository's reflog and on the remote
server. A reader of the eventual publication can verify the lock
state by checking out the lock tag directly and inspecting the
artifacts as they existed at lock time.

\subsection{Execution and operational-fix protocol}
\label{subsec:execution}

The locked methodology is executed once on the evaluation data. The
execution produces a primary output set: numeric results, log files,
and any intermediate artifacts specified in the pre-registration.
These outputs are committed to the repository alongside the analysis
code and form the basis for the RESULT document. Re-executions are
permitted only for reproducibility verification: re-running the
locked code on the locked data to confirm that recorded outputs
match. Re-executions do not produce new RESULT documents.

Bugs in the analysis code discovered after the lock commit may be
corrected, but only under constrained conditions. An
\emph{operational fix} is permitted if and only if (a) the fix does
not alter the analysis specification as recorded in the
pre-registration, (b) the fix is committed as an explicit commit
with a descriptive message that references both the pre-registration
and the bug being addressed, (c) the fix is documented in the RESULT
document with the commit reference recorded, and (d) the execution
that produced the recorded outputs used the fixed code rather than
the original buggy code. A change that alters a multiple-testing
correction method, a significance threshold, or a decision rule is
not an operational fix; it requires re-locking with a new
pre-registration version, a new cross-model review, and a new lock
tag.

The ERSAF v1 execution required three operational fixes: an
encoding configuration for non-ASCII characters in log output, a
standard-output buffering directive for progress reporting on the
Windows host, and an attribute-initialization fix in a change-point
detector cost function (commit \texttt{552cc0c}). Each is documented
in the disclosures section of the RESULT document. Where the
distinction between an operational fix and a specification change
is ambiguous, the framework errs on the side of re-locking.

A failure mode worth naming explicitly: an unscrupulous operator
could in principle execute the locked analysis, observe an
unfavorable result, declare the run ``buggy,'' commit an
``operational fix'' that subtly biases the analysis toward a more
favorable outcome, and re-execute. The framework's protection
against this failure mode does not come from the operational-fix
protocol's text alone but from the artifact chain the protocol
generates. The original ``buggy'' execution remains in the
repository's commit history; the operational-fix commit is itself
visible alongside it; the cross-model reviewer (whose written
review is committed at lock time) can be invited to dispute the
classification of a post-lock change as ``operational'' versus
``specification''; and any reader of the eventual publication can
inspect the commit chain between the lock tag and the result tag
to look for patterns---successive operational fixes moving the
result in a single direction, in particular---that warrant
challenge. The framework's claim is that an AEPF artifact chain
makes this failure mode harder to commit without leaving evidence,
not that it makes the failure mode impossible.

\subsection{RESULT document}
\label{subsec:result-document}

After execution, the operator produces the RESULT document. The
RESULT contains a headline verdict, one of POSITIVE, NULL, or
AMBIGUOUS as defined by the pre-registered decision rules; per-baseline,
per-observable statistical findings, including all test statistics
and any confidence or credible intervals; required disclosures,
covering pre-registered limitations, systematics that became apparent
during execution, and any operational fixes applied; and a
\emph{constrained-interpretation section} that, when multiple
permissible interpretations apply to the result, documents each frame
side-by-side without claiming one stronger than the other given the
locked methodology's scope. The RESULT also includes a forward-looking
section describing what the methodology deliberately did not test,
providing context for follow-up work and bounding the claims the
current result can support.

The RESULT is committed to the repository and tagged following the
convention \texttt{<framework>-v<version>-result}
(e.g., \texttt{ersaf-v1.0-result}), and the result tag is pushed to
the same public remote as the lock tag. For the ERSAF audit, the
headline verdict was NULL: no Bonferroni-corrected detection survived
in any baseline model. A single Benjamini--Hochberg false discovery
rate firing at $z \in [0.024, 0.038]$ was reported under two
interpretive frames: a formal frame, in which the firing is
consistent with smooth dark-energy absorption at the chosen
significance level, and a systematic frame, in which the firing is
consistent with the low-$z$ peculiar-velocity systematic disclosed
in the pre-registration. Neither frame was claimed stronger than the
other; the RESULT document records both \cite{dixon2026ersaf}.

\subsection{Calibration assessment as standing requirement}
\label{subsec:calibration-req}

In domains where the evaluation outputs probability estimates --
classifier confidence scores, forecast probabilities, regression
uncertainty estimates -- AEPF requires calibration assessment
alongside whichever discriminative metrics are appropriate to the
domain. The calibration-assessment requirements specify the Brier
score \cite{brier1950} and the Expected Calibration Error
\cite{naeini2015ece, guo2017calibration} as minima, accompanied by a
reliability diagram. The requirement is treated as a standing
requirement: a pre-registration that does not specify calibration
assessment for a probability-output evaluation does not pass
cross-model review.

The motivation is that strong discriminative metrics in the presence
of poor calibration produce systems that are confidently wrong --- a
failure mode that does not appear in accuracy-only or AUROC-only
evaluation \cite{guo2017calibration, minderer2021revisiting}.
Calibration audits performed under AEPF discipline have documented
this failure mode empirically in multiple transformer-based language
models \cite{dixon2026distilbert, dixon2026toxicbert,
dixon2026chatgptdetector}, in each case revealing systematic
overconfidence that would have been invisible under accuracy-driven
evaluation alone.

\subsection{Instrument validation of the calibration audit}
\label{subsec:instrument-validation}

The standing calibration requirement
(Section~\ref{subsec:calibration-req}) presumes the calibration-audit
rule is itself a sound instrument: that it accepts genuinely
calibrated probabilities and rejects miscalibrated ones. Because the
framework leans on that rule to surface a failure mode invisible to
accuracy-only evaluation, the rule was validated as an instrument
against synthetic ground truth---a methods-level exercise, distinct
from a locked AEPF study, that characterizes the \emph{tool} rather
than making a finding about any model or dataset. The audit rule
under test was the one used verbatim in the transformer calibration
audits (Section~\ref{subsec:calibration-case}): ten equal-width bins,
a bin passing only if adequately populated and its mean predicted
probability falls inside the Wilson interval of its observed
frequency, scored against a base-rate climatology. The validation ran
the rule over 100 replications on each of several synthetic arms whose
calibration was known by construction, under a single master seed.

\paragraph{Operating characteristics.}
The results separate cleanly into two distinct error properties.
\emph{Soundness} (Type~I, the false-accept of broken calibration) was
zero across all four miscalibration families tested---overconfident,
underconfident, positive-shift, and negative-shift logit
perturbations were each rejected in 100 of 100 replications. Not a
single replication of broken calibration was accepted. \emph{Power}
(Type~II, the false-reject of genuine calibration) was zero under
well-spread calibrated predictions---uniform and bimodal arms were
each accepted 100 of 100---and approximately 11\% under a
heavily-skewed-but-genuinely-calibrated $\mathrm{Beta}(2,5)$ arm,
which concentrates predictions in the low range and leaves the upper
bins sparse. Noise tolerance was monotone: the instrument accepted
small estimation noise and rejected large noise in the predicted
direction.

\paragraph{The mechanical verdict and what it rediscovered.}
The validation carried a single pre-committed success bar: every
positive-control arm classified as calibrated in at least 90 of 100
replications, and every negative-control arm rejected in at least 90
of 100. Two of the three positive arms cleared it outright; the
skewed $\mathrm{Beta}(2,5)$ arm came in at 89 of 100---one replication
under the line. The conjunction therefore failed, and the verdict is
reported as FAIL: no kinder seed, no softened threshold, no
relabeling. The skewed arm's $\sim$11\% false-reject rate is not the
instrument failing to recognize calibration---the outcomes were drawn
from the predictions, so there is no miscalibration in the arm to
recognize---but the empirical signature of a power limitation the
audit template had named \emph{a priori}: highly skewed distributions
reduce the bin-wise statistical power of the calibration test. An
a-priori caveat thus acquired a measured operating characteristic.

\paragraph{The design lesson.}
The single pre-committed bar conflated two distinguishable properties:
instrument soundness (which held completely) and power under skew (on
which the skewed arm fell one replication short). A correctly
stratified pre-commitment would have made soundness and a
well-spread power baseline binding while reporting power under a
documented difficulty---skew---as characterization rather than as a
gate, because the template had predicted reduced power there in
advance. This is recorded as the framework's v1.1 refinement for
instrument-validation studies: pre-committed criteria must be
stratified by the property under characterization rather than
collapsed into a single conjunction across heterogeneous arms.

\paragraph{A contained claim, not a general principle.}
The disciplined report of this near-miss FAIL strengthens the
credibility of the instrument-validation finding---the failure was
small, the diagnosis was sharp, and the substantive interpretation
was constrained by an a-priori template note rather than improvised
after the fact. This claim is deliberately contained to this case. It
is \emph{not} a general principle that a FAIL outranks a PASS, and it
must not be read as one. If AEPF reports routinely failed their
pre-committed bars and leaned on substantive interpretation to
recover, pre-committed bars would become decoration and the framework
would invert into the very practice it exists to prevent. The
near-miss strengthens credibility here because it is bounded by a
small failure, a sharp diagnosis, and an a-priori commitment; it
earns its containment case-by-case and claims nothing beyond this
case.

\subsection{Null-result publication parity}
\label{subsec:null-parity}

AEPF requires NULL and POSITIVE results to be published with
identical artifact discipline. The pre-registration's publication
commitment binds the operator to publish the RESULT regardless of
outcome direction. A NULL verdict is not a non-event: it is the
project's output, produced under the same locked methodology that
would have produced a POSITIVE verdict had the data supported one,
and reported with the same level of statistical detail and the same
artifact preservation requirements.

The framework's null-result publication discipline addresses
publication bias \cite{ioannidis2005} at the methodology-lock stage
rather than at the journal-submission stage. By the time the result
is observed, the publication commitment is already in place; the
operator cannot retract a study they intended to publish only if it
produced a positive finding. The ERSAF NULL verdict and its
accompanying RESULT document, published under the lock-tag
discipline described above, illustrate the operational consequence:
a null finding is treated as a substantive contribution to the
literature rather than as a failed project.

\subsection{Viability gate (a refinement under observation)}
\label{subsec:viability-gate}

A pattern that emerged in operational practice during the
preparation of one of the case studies in
Section~\ref{sec:case-studies} extends the framework backward to
the pre-specification phase. When a candidate pre-registration is
drafted whose informational yield depends on an empirical
assumption that has not been checked---particularly an
independence assumption between predictors, or a structural prior
that makes the answer predictable in advance---AEPF as articulated
in Sections~\ref{subsec:prereg-review} through
\ref{subsec:null-parity} has no explicit handle for the question
``is this study worth locking at all?'' The framework's entry
point assumes the question is already specifiable enough to
pre-register; it does not address the case where the question is
specifiable but the answer is structurally over-determined before
data are observed.

The \emph{viability gate} is a structured exploratory diagnostic,
run before lock, with a pre-committed decision rule about whether
to proceed with the candidate pre-registration. The diagnostic is
explicitly non-confirmatory: it cannot be cited as confirmatory
evidence in any subsequent published audit, and any findings that
prove substantively interesting in their own right must be
re-derived under proper lock discipline before publication. Its
purpose is exclusively to gate the decision to lock, not to
substitute for the locked study.

The discipline that distinguishes a legitimate viability gate from
an exploratory peek that would contaminate the subsequent
pre-registration is a constraint on what the diagnostic is
permitted to inform. The pre-committed decision rule must be about
\emph{viability}---does the candidate study have structural room
to return an informative answer?---and never about
\emph{methodology selection}---which combination rule, threshold,
or transformation should the pre-registration lock? The moment a
viability gate begins answering methodology-selection questions,
it has become an unblinded peek and the candidate
pre-registration's methodology has been data-informed in the way
AEPF is designed to prevent.

Trigger conditions for performing a viability gate, drawn from
observed practice: independence assumptions are critical to the
study's informational content; structural priors suggest a NULL
result for reasons unrelated to the substantive question;
sample-size feasibility is genuinely uncertain; or the candidate
pre-registration has a methodologically defensible structure but
its informational yield depends on an empirical assumption that
has not been checked. The Dynametrix-HRRR case study
(Section~\ref{subsec:dynametrix-hrrr-case}) is the worked example.

A note on the edge between viability and methodology selection.
The Dynametrix-HRRR diagnostic uses a trichotomous decision rule:
PROCEED, PROCEED\_WITH\_DISCLOSURE (which adds a required
disclosure to the candidate pre-registration's required-disclosures
section if triggered), and DO\_NOT\_PROCEED. The middle branch
sits at the edge of the viability/methodology-selection
distinction. The framework's position is that a required disclosure
added to the pre-registration is an interpretive constraint rather
than a methodology change: it does not alter what tests are run,
which corrections are applied, or how the verdict is computed; it
constrains how a positive verdict, if produced, may be reported.
On this reading the middle branch remains viability-side. The case
is the first operational application of the gate, and the
boundary between ``interpretive constraint addition'' and
``methodology change'' is a frontier the gate's continued use will
need to map empirically before formalization.

This subsection is included as a refinement under observation
rather than as a fully formalized v1.0 component of AEPF. The
pattern has, at the time of writing, one operational application.
Further applications across heterogeneous contexts will determine
whether the trigger conditions and constraints as stated here hold
up, or whether they require revision before formalization in a
subsequent version of the framework.

\section{Case Studies}
\label{sec:case-studies}

This section presents four case studies applying AEPF discipline
to substantively different evaluation contexts. The first
(Section~\ref{subsec:ersaf-case}) is a cosmological audit testing
for non-smooth residual structure in joint expansion-history
residuals---a full-lifecycle study that locked, executed, and
reported a NULL verdict. The second
(Section~\ref{subsec:calibration-case}) is a series of calibration
audits performed on transformer-based language models released to
public model registries---a demonstration of the framework's
portability across multiple model evaluations under shared
discipline. The third
(Section~\ref{subsec:rc-case}) is the Relational Coherence audit
(RC v1)---a pre-registered framing-robustness test of an open-weight
language model that locked, executed, and reported a NULL verdict on
a single commodity GPU. The fourth
(Section~\ref{subsec:dynametrix-hrrr-case}) is the
Dynametrix-HRRR independence diagnostic---the first operational
application of the viability gate pattern described in
Section~\ref{subsec:viability-gate}, illustrating the framework's
handling of pre-specification dispositions. The four case studies
share no domain-specific machinery; all four share the AEPF
artifact discipline. Their juxtaposition demonstrates the
framework's portability across heterogeneous evaluation contexts
and shows how the same operational artifacts---pre-registration,
materialization manifest, lock-time commitment, single-execution
discipline, RESULT document with constrained interpretation, and
where applicable the viability gate---can be applied without
modification to problem domains that share no other methodological
vocabulary.

\subsection{ERSAF: a cosmological residual-structure audit}
\label{subsec:ersaf-case}

The Expansion Residual Structure Assessment Framework (ERSAF)
\cite{dixon2026ersaf} applied AEPF discipline to a search for
non-smooth residual structure in the joint expansion-history
residuals of the Pantheon+ Type~Ia supernova compilation
\cite{brout2022pantheon} and the DESI Data Release~1 Baryon Acoustic
Oscillation measurements \cite{adame2024desi3, adame2024desi4},
relative to three locked smooth dark-energy baselines: flat \LCDM{},
flat \wCDM{}, and flat \wowaCDM{}.

\paragraph{Pre-registration and cross-model review.}
The ERSAF pre-registration document was finalized prior to any
analysis touching the evaluation data and specified the full
analytical pipeline: baseline cosmology evaluation against the
locked datasets, residual computation, two test families (a
continuous wavelet transform with Torrence \& Compo cone-of-influence
masking and a PELT change-point detector with weighted Gaussian
cost), multiple-testing correction with Bonferroni family-wise error
rate as primary and Benjamini--Hochberg false discovery rate as
secondary, and decision rules classifying outcomes as POSITIVE,
NULL, or AMBIGUOUS at pre-specified significance thresholds. The
pre-registration also disclosed methodological limitations,
including a low-redshift peculiar-velocity systematic acknowledged
prospectively as a known blind spot of the chosen residual
construction.

The pre-registration underwent two rounds of cross-model review
before lock. The first round produced numbered findings classified
by severity; the operator addressed each finding and committed the
responses alongside the pre-registration. A second focused review
examined only the analytical changes introduced by the first
round's remediation, ensuring those changes had not introduced new
methodological issues. Both review records were committed to the
repository as part of the lock commit.

\paragraph{Materialization and lock.}
The materialization manifest recorded the SHA-256 hashes of three
classes of artifact: the Pantheon+ data and covariance files as
released by the Pantheon+ collaboration; the DESI DR1 BAO
12-observable values, materialized from the upstream
\texttt{CobayaSampler/bao\_data} repository at a named commit, with
each observable's value, redshift, and bin label cross-checked
against the published Tables of \cite{adame2024desi3,
adame2024desi4} within the half-unit tolerance of the published
decimal precision; and the analysis script itself, with its
SHA-256 recorded in the pre-registration document. All artifacts
were committed in a single lock commit and tagged
\texttt{ersaf-v1.0-lock}; the tag was pushed to the public remote
before unblinding.

\paragraph{Execution and operational fixes.}
The locked methodology was executed once. The execution required
three operational fixes, each committed as a distinct repository
commit and documented in the disclosures section of the RESULT
document: a UTF-8 encoding configuration for non-ASCII characters
in log output; a standard-output buffering directive
(\texttt{python -u}) for progress reporting on the Windows host
where the run was performed; and an attribute-initialization fix
in the PELT cost function (\texttt{self.signal = signal} added to
the fit method), committed as \texttt{552cc0c}. None of the three
fixes altered the analysis specification; each addressed an
implementation defect in the locked code rather than a
methodological decision. The execution that produced the recorded
outputs used the fixed code.

\paragraph{RESULT and interpretation.}
The headline verdict was NULL: no Bonferroni-corrected detection
survived in any of the three baseline models. A single
Benjamini--Hochberg false discovery rate firing was observed in the
wavelet test family at $z \in [0.024, 0.038]$, corresponding to a
single coefficient at wavelet scale 1.2854 bins. The RESULT
document reports this firing under two interpretive frames,
documented side-by-side without claiming one stronger than the
other given the locked methodology's scope.

The \emph{formal frame} interprets the BH FDR firing as a
statistically marginal feature, consistent with smooth dark-energy
absorption at the chosen significance level: the absence of any
Bonferroni-firing detection indicates that the joint residuals
contain no structure incompatible with the smooth baselines once
appropriate multiple-testing correction is applied.

The \emph{systematic frame} interprets the same firing as the
expected signature of the low-$z$ peculiar-velocity systematic
disclosed prospectively in the pre-registration. The firing's
location at $z = 0.024$--$0.038$ falls within the redshift range
where peculiar-velocity contamination is known to be most
significant relative to cosmological signal, and the
single-coefficient pattern is consistent with a localized
systematic rather than a broadband cosmological feature.

The RESULT document explicitly declines to elevate either frame
over the other on the basis of the locked methodology alone.
Resolution between frames would require a follow-up analysis
specifically targeting the low-$z$ peculiar-velocity systematic
under a separately locked methodology---work that ERSAF v1
deliberately did not perform and that the RESULT's forward-looking
section identifies as the natural successor.

The RESULT was committed and tagged \texttt{ersaf-v1.0-result};
the tag was pushed to the public remote alongside the lock tag.

\paragraph{What ERSAF demonstrates about AEPF.}
The ERSAF case demonstrates AEPF's behavior under a NULL outcome.
A non-AEPF audit observing a single BH FDR firing at low redshift
could plausibly have foregrounded that firing as a marginal
positive result, with the systematic interpretation relegated to
supplementary material. Under AEPF discipline, the publication
commitment locked at the pre-registration stage required the NULL
verdict be reported with the same artifact discipline as a
POSITIVE result would have received, and the constrained-interpretation
section required the systematic frame be reported with equal
standing to the formal frame. The reader of the RESULT receives
both interpretations, the methodology's scope is explicitly
bounded, and the natural follow-up work is identified rather than
implied.

\subsection{Transformer calibration audits}
\label{subsec:calibration-case}

A second set of AEPF case studies applied the framework to
calibration assessment of transformer-based language models
released to public model registries. Three audits were performed
under separate pre-registrations and separate lock tags, each
following the AEPF lifecycle: DistilBERT fine-tuned on the
Stanford Sentiment Treebank binary task (DistilBERT-SST2)
\cite{dixon2026distilbert}; Toxic-BERT, a multi-label toxicity
classifier \cite{dixon2026toxicbert}; and a RoBERTa-based
ChatGPT-detection model \cite{dixon2026chatgptdetector}.

\paragraph{Common framework, separate locks.}
Each audit produced its own pre-registration document, its own
cross-model review, its own materialization manifest pinning both
the evaluation corpus and the model revision identifier, its own
lock commit and lock tag, and its own RESULT document. The three
audits did not share lock tags, did not share materialization
manifests, and were not pooled at the analysis stage. Sharing the
framework did not mean sharing the artifact chain: each model and
each evaluation context received an independent AEPF lifecycle.

\paragraph{Calibration as primary deliverable.}
Each pre-registration specified, per Section~\ref{subsec:calibration-req}
of this paper, the Brier score \cite{brier1950} and Expected
Calibration Error \cite{naeini2015ece, guo2017calibration} as
primary deliverables, accompanied by reliability diagrams. The
discriminative metrics appropriate to each task---accuracy and
F1 for binary classification, per-label AUROC for the multi-label
toxicity task, accuracy and discrimination-threshold sensitivity
for the ChatGPT detector---were reported as secondary deliverables.
This ordering is the operational consequence of
Section~\ref{subsec:calibration-req}: a model's probability
estimates must be assessed for reliability before its
discriminative outputs are characterized.

\paragraph{Materialization of model and corpus.}
The materialization manifests for the three audits each pinned the
evaluation corpus by file SHA-256 (the SST-2 development set for
DistilBERT-SST2, a stratified subsample of the toxicity evaluation
data for Toxic-BERT, and a curated set of human-written and
machine-generated text for the ChatGPT detector) and pinned the
model itself by repository revision identifier. The revision pin
is essential because models in public registries are mutable
artifacts: a model identified only by repository name may be
silently updated by the maintainer between materialization and
re-execution, breaking reproducibility. Pinning the revision
identifier in the manifest ensures that a re-execution loads
exactly the bytes that were evaluated at lock time.

\paragraph{Findings and pattern.}
Each of the three audits documented systematic miscalibration in
the model under evaluation. Table~\ref{tab:calibration-summary}
summarizes the headline metrics. The reliability diagrams
\cite{dixon2026distilbert, dixon2026toxicbert,
dixon2026chatgptdetector} display overconfidence in the
high-probability bins for the binary and multi-label
classification audits, and a more complex calibration profile for
the ChatGPT detector reflecting the task's open-ended distribution
shift. In all three cases, the discriminative metrics taken in
isolation would have suggested adequate model performance: accuracy
ranged from 67.8\% (ChatGPT detector) to 94.5\% (Toxic-BERT). The
calibration assessment revealed that probability estimates from
these models, treated as reliable confidence outputs, would
systematically mislead downstream consumers. The ChatGPT detector
case is the most striking: a Brier Skill Score of $-0.16$ against a
0.5 base-rate climatology means the model's probability outputs are
\emph{worse} than the trivial reference forecast that always
predicts the base rate---a finding invisible to the model's
67.8\% accuracy score taken alone.

\begin{table}[h]
\centering
\caption{Headline calibration metrics across the three transformer
calibration audits. Brier = Brier score (lower is better); BSS =
Brier Skill Score against base-rate climatology (positive is better,
negative is worse than reference); ECE = Expected Calibration Error.
Outcome category per the audit's pre-registered decision rule.}
\label{tab:calibration-summary}
\begin{tabular}{lrrrrrl}
\toprule
Model & $N$ & Accuracy & Brier & BSS & ECE & Outcome \\
\midrule
DistilBERT-SST2 & 872 & 0.911 & 0.083 & $+0.666$ & 0.072 & Not calibrated \\
Toxic-BERT & 5{,}000 & 0.945 & 0.044 & $+0.349$ & 0.022 & Drift detected \\
ChatGPT-detector & 2{,}000 & 0.678 & 0.291 & $-0.163$ & 0.279 & Not calibrated \\
\bottomrule
\end{tabular}
\end{table}

\paragraph{What the calibration audits demonstrate about AEPF.}
The three calibration audits demonstrate AEPF's domain portability:
the same framework that locked a cosmological audit of dark-energy
extensions locked three audits of transformer probability outputs
without modification. They also demonstrate the operational value
of treating calibration as a standing requirement
(Section~\ref{subsec:calibration-req}). In the absence of that
requirement, calibration assessment would have been a project-level
decision, and the three audits could plausibly have reported
discriminative metrics alone---missing the failure mode the audits
were specifically designed to detect. The standing requirement
forecloses that failure path at the pre-registration stage.

\subsection{Relational Coherence v1: a pre-registered open-weight LLM null}
\label{subsec:rc-case}

A third full-lifecycle case study applied AEPF discipline to an
open-weight language model on commodity hardware, returning a clean
NULL. The Relational Coherence audit (RC v1)
\cite{dixon2026rcv1} pre-registered a single gated question: does the
\emph{behavioral-equivalence graph topology} of a model's outputs
under semantically-equivalent framing variation reveal consistency
structure that an aggregate pairwise-agreement rate conceals?

\paragraph{Question and pre-committed test.}
For each query, the model was run on $N_{\mathrm{var}} = 50$
framing-and-layout variants that preserve the verbatim question and
option text and vary only the instructional wrapper and option
layout. From the resulting output graph two scalars were defined: an
aggregate pairwise consistency rate $C$ (edge density) and a
pair-fraction-in-same-component $F$ (the probability that two
randomly chosen variant outputs share a behavioral component). The
two are not algebraically reducible to one another---graphs of equal
edge density can differ in connectivity---and the pre-registration
gated on whether topology $F$ separates queries that $C$ rates as
equivalent. The pre-committed decision rule: among comparable
query-pairs ($|\Delta C| \leq 0.05$), at least $k = 0.20$ must show
topological separation ($|\Delta F| \geq 0.10$). The pre-registration
recorded prospectively that a high-framing-robustness regime, in
which $F$ saturates near one across queries, would yield a NULL, and
committed to publishing that NULL under the same discipline as a
positive result.

\paragraph{Substrate, model, and lock.}
The substrate was the canonical MMLU release pinned by dataset
revision: 300 queries balanced across three subjects
(high-school mathematics, professional law, miscellaneous),
selected deterministically by content hash. The model under test was
Qwen2.5-7B-Instruct-AWQ, a 4-bit quantized open-weight
instruction-tuned model pinned by repository revision and served via
vLLM under greedy decoding. The edge predicate used a pinned
sentence-embedding model, with the equivalence threshold $\tau$
calibrated to a held-out, disjoint slice and frozen before any
main-run inference under a separate \texttt{rc-v1-calibrated} tag.
The pre-registration, analysis script, and materialization manifest
were committed in a single lock commit and tagged
\texttt{rc-v1-lock}; the materialization manifest pinned the dataset,
the model weights, and the embedding model each by revision, so that
a re-execution loads exactly the bytes evaluated at lock time.

\paragraph{Result and verdict.}
The frozen threshold was $\tau = 0.8017$. Across 300 queries (none
dropped for inference failure), $21{,}946$ comparable query-pairs
were identified; $148$ showed topological separation, a proportion
of $0.0067$ against the pre-committed $k = 0.20$. The headline
verdict was NULL. The null was driven by $F$-saturation: mean
$F = 0.993$ (minimum $0.65$), against a more variable $C$
(mean $0.930$, minimum $0.296$). The secondary, ungated
characterization---the count of high-accuracy, low-$F$ queries that
would signal cross-framing fragmentation despite correct label
recovery---was zero. The behavioral graphs were, in short, almost
never fragmented: the model is highly framing-robust on this task
family, and the topology $F$ found no consistency structure for $C$
to have concealed because there was little structure to find. This
is the high-robustness regime the pre-registration named in advance
as a NULL-producing outcome, now observed.

\paragraph{What RC v1 demonstrates about AEPF.}
The RC v1 case demonstrates the framework producing a clean null
where the structural prior favored one, published with full
positive-parity rather than abandoned. Because the high-robustness
NULL was named prospectively in the pre-registration, the verdict is
not a disappointed positive recast as a null; it is the
pre-specified outcome of a single locked execution. The case also
exercises the full lock${\rightarrow}$calibrate${\rightarrow}$result
single-execution pipeline---a separately tagged calibration step
freezing a held-out threshold, followed by a single main run---on a
quantized open-weight model running on a single commodity GPU,
demonstrating that the artifact discipline AEPF prescribes carries
to small-scale, locally reproducible language-model evaluation
without modification.

\subsection{Dynametrix-HRRR: a viability-gate case}
\label{subsec:dynametrix-hrrr-case}

The fourth case study applies AEPF discipline to a question whose
answer turned out to be structurally over-determined before the
candidate study could be locked. The candidate question:
\emph{does combining Dynametrix v3 calibrated severe-weather
probabilities with NWS HRRR-derived probabilities produce a
measurably better-calibrated joint prediction than either source
alone, over a defined location set and forecast horizon?}

\paragraph{Candidate pre-registration.}
A draft pre-registration was prepared specifying Brier Skill Score
of the combined predictor against HRRR alone as the primary
metric, with simple arithmetic averaging as the primary
combination rule, three secondary combination rules subject to
Bonferroni correction, and a publication commitment binding the
operator to publish whatever verdict the data returned. The draft
was structured to be lockable: data sources were named, the
forecast horizon was fixed, the evaluation window was defined,
and required disclosures including an independence-assumption
caveat were prospectively stated.

\paragraph{Pre-lock viability diagnostic.}
Before advancing the draft to cross-model review and lock, an
exploratory independence diagnostic was performed under an
explicitly non-confirmatory protocol with a pre-committed decision
rule. The diagnostic measured the Pearson correlation of residual
errors between Dynametrix v3 and a provisional HRRR-derived
predictor over a defined past-data window
(2026-05-18 through 2026-05-24, 2{,}016 location-hours,
375 verified severe-weather events). Decision thresholds: residual
correlation $r < 0.4$ proceeds unconstrained; $0.4 \leq r \leq 0.7$
proceeds with a required disclosure of bounded independence;
$r > 0.7$ does not proceed.

\paragraph{Diagnostic outcome.}
The measured Pearson correlation was $r = 0.9839$ (95\% bootstrap
CI $[0.9816, 0.9860]$, $B = 1{,}000$ stratified resamples). Per-category
correlations: tornado $r = 0.949$ ($n = 1{,}710$, 69 events), hail
$r = 0.944$ ($n = 1{,}706$, 65 events), wind $r = 0.979$
($n = 1{,}882$, 241 events). Mean residuals indicated systematic
miscalibration of both predictors in opposite directions: mean
Dynametrix residual $-0.18$ (systematic over-prediction relative to
an 18.6\% observed event rate), mean HRRR-derived residual $+0.18$
(systematic under-prediction under the provisional transformation).
Pearson correlation is invariant to such marginal shifts, so the
headline correlation reflects signal-tracking agreement between
the two predictors despite their disagreement on absolute
probability calibration. The verdict from the pre-committed
decision rule: \texttt{DO\_NOT\_PROCEED}.

\paragraph{Disposition.}
The candidate pre-registration was archived as ``considered but
deferred'' rather than advanced to cross-model review and lock
\cite{dixon2026dynametrixhrrr}. The archived artifact preserves
the full draft pre-registration alongside a header documenting
the diagnostic verdict that retired it, and identifies the
conditions under which the question might be revisited: a
meaningfully different HRRR-to-probability transformation, a
revised Dynametrix feature set with non-overlapping upstream
signal, or substitution of HRRR with a less-correlated reference
forecast. Absent one of these conditions, the question is closed.

\paragraph{What the Dynametrix-HRRR case demonstrates about AEPF.}
The case demonstrates the framework's handling of an outcome
neither full-lifecycle case study produces: a pre-specification
disposition in which the candidate study is retired before lock
because its informational yield is structurally constrained.
Without a viability-gate pattern, the candidate study could
plausibly have advanced to lock and execution, consumed analyst
time on a methodologically clean run, and produced a NULL or
inconclusive RESULT whose interpretive value would have been
limited by the structural prior. With the viability gate, the
structural prior was checked first, the candidate disposition was
documented, and the analyst time freed for a different audit
target. The exploratory-not-confirmatory framing of the gate
ensures the disposition itself is not laundered into a
confirmatory finding: the diagnostic record explicitly cannot be
cited as evidence about independence between Dynametrix v3 and
HRRR more broadly, only as the gate that retired this specific
candidate study.

\paragraph{Joint observation across the case studies.}
The four case studies together illustrate four complementary
patterns under AEPF discipline. The ERSAF case shows the
framework's behavior when a locked full-lifecycle study returns a
NULL verdict: the result is published with the same discipline as
a positive verdict would have received, the dual interpretive
frames are recorded without one being promoted above the other,
and the deliberate non-coverage of the methodology is documented
in the forward-looking section. The transformer calibration
audits show the framework's behavior across multiple evaluations
in a shared domain: each lock chain is independent, each RESULT
stands on its own artifact chain, and the standing calibration
requirement ensures that a failure mode (overconfident probability
outputs) that is invisible under accuracy-only evaluation is
surfaced across all three audits. The Relational Coherence audit
shows the framework's behavior on an open-weight language model under
greedy decoding: a high-framing-robustness regime named prospectively
in the pre-registration as NULL-producing returned exactly that NULL,
through the full lock--calibrate--result single-execution pipeline on
a single commodity GPU. The Dynametrix-HRRR case shows
the framework's behavior when a candidate study is structurally
not worth locking: the viability gate produced a documented
pre-specification disposition, retiring the candidate
pre-registration before it consumed lock-chain resources. None of
the four patterns depends on the domain; all four depend on the
artifact discipline AEPF prescribes.

\section{Discussion}
\label{sec:discussion}

The claim AEPF makes is narrower than may be apparent from its
artifact discipline. The framework does not eliminate analyst freedom,
does not guarantee correct conclusions, and does not replace the
domain-specific judgment that empirical evaluation requires. Its claim
is that analyst freedom, where exercised, becomes visible and
cryptographically dated; that the methodology under which a result was
produced is independently verifiable against a fixed reference point;
and that the institutional commitment to publish a result has been
recorded before the result was observed. What changes under AEPF is
not the rate at which evaluations produce correct conclusions but the
rate at which a reader can confidently assess what the conclusions
mean.

\paragraph{The burden-of-proof inversion.}
In conventional evaluation, the burden of demonstrating that a result
is robust falls on the reader. The reader must assess, often without
access to the underlying artifacts, whether the methodology was
specified before unblinding, whether the data sources reported are
the data sources actually used, whether the statistical procedures
were chosen post-hoc to surface the reported finding, and whether
alternative methodological choices would have produced different
conclusions. AEPF inverts this burden: a reader of an AEPF-locked
study can verify these properties directly against the lock-tag
artifacts. The reader's role becomes one of confirming that the
pre-registered methodology matches the executed methodology, that
the materialization manifest matches the cited data sources, and
that any operational fixes documented in the RESULT do not cross
into specification changes. Verification replaces inference.

\paragraph{The role of peer review.}
AEPF does not replace peer review. The pre-registration document
addresses methodology decisions made before unblinding; peer review
at publication addresses the broader questions of whether the
methodology was well-chosen, whether the evaluation question was
worth asking, and whether the conclusions drawn from the result are
adequately bounded by the evidence. These are distinct functions.
Cross-model review under AEPF is a methodology-stage review, intended
to catch avoidable methodological errors before they are locked into
the artifact chain; peer review at publication is a results-stage
review, intended to situate the locked result within the relevant
literature and assess its broader claims. AEPF strengthens the
artifact chain peer review operates on; it does not substitute for
the substantive judgment peer review provides.

\paragraph{Null verdicts as substantive contributions.}
Perhaps the most consequential operational change AEPF introduces is
its treatment of NULL verdicts. Under conventional publication
incentives, a project that does not produce a positive finding tends
either to remain unpublished or to be reported with substantially
less discipline than a positive finding would have received
\cite{ioannidis2005, nosek2015open}. AEPF's publication commitment,
locked at the pre-registration stage, removes the decision of whether
to publish from the post-unblinding analytic phase. The operator who
locks a pre-registration is institutionally bound to publish whatever
the data return. The ERSAF NULL verdict, published under this
discipline, is a worked illustration: a result that under conventional
incentives might have been characterized as ``inconclusive'' or
``promising but inconclusive'' is instead a substantive contribution,
with the dual interpretive frames documented at the same level of
detail as a positive finding would have warranted.

\paragraph{Instrument soundness as a precondition for null evidence.}
A null verdict is only evidence about the world if the instrument
that produced it is sound. The validation of the calibration audit
(Section~\ref{subsec:instrument-validation}) establishes that
soundness directly---zero false-accept of broken calibration across
the tested families---and in doing so surfaced a refinement to the
framework's pre-commitment discipline. When the object under test is
an evaluation instrument validated against synthetic ground truth
rather than a hypothesis about the world, the pre-committed success
criteria must be stratified by the property under
characterization---soundness, power baseline, and power under named
difficulties---rather than collapsed into a single conjunction across
heterogeneous arms. A single bar spanning arms that probe different
properties can fail for reasons that do not bear on the property that
matters, obscuring the instrument's operating characteristics rather
than revealing them. This v1.1 refinement is contained to
instrument-validation studies; it does not weaken the pre-committed
bar for hypothesis-testing studies, where a single decision rule
remains appropriate.

\paragraph{What AEPF does not change.}
AEPF does not change what is true. A well-locked methodology
applied to inadequate data will produce a well-locked result that
is uninformative about the underlying question. A pre-registration
that asks the wrong question will produce an answer to the wrong
question, regardless of the lock tag's cryptographic strength.
AEPF's value is procedural: it constrains the relationship between
the question asked, the methodology applied, and the verdict
reported. The substantive correctness of the resulting work
depends on judgments AEPF cannot enforce---whether the question is
worth asking, whether the methodology is appropriate to it,
whether the result's conclusions are bounded by what the locked
artifacts can support. These remain the responsibility of the
operator, the cross-model reviewer, and the publication reviewer.

\section{Limitations}
\label{sec:limitations}

AEPF is a workflow discipline rather than a statistical guarantee.
Five limitations bound its claims, each of which the framework
acknowledges rather than purports to resolve.

\paragraph{Retrospective pre-registration.}
AEPF cannot detect a researcher who performs unblinded exploratory
analysis privately and only then writes a pre-registration that
matches the exploratory finding. The cryptographic discipline
attests that the pre-registration was committed at a particular
moment in time; it does not attest that no prior unblinded
exploration occurred. Mitigation depends on social and institutional
factors outside the framework: declared conflicts of interest, the
operator's incentives, and the willingness of cross-model reviewers
to flag pre-registrations that appear unusually well-fitted to a
specific hypothesized outcome. The framework's response to this
limitation is to make the artifact chain a credible negative
signal---a researcher who has performed retrospective pre-registration
takes on the burden of explaining why their pre-registration anticipated
the eventual outcome with such precision---rather than to claim a
solution.

\paragraph{Cryptographic infrastructure assumptions.}
The framework's verifiability properties rely on the integrity of
the version-control infrastructure and the public availability of
the remote where lock tags are pushed. A repository whose remote
ceases to be available, whose history is rewritten, or whose tag
references are broken loses the artifact chain's verifiability. The
mitigation is mirroring lock tags to multiple independent remotes
and, where the stakes warrant, depositing the artifact chain with
an institutional preservation service. These are operational
practices the framework recommends but cannot enforce
cryptographically.

\paragraph{The operational-fix boundary.}
The distinction between an operational fix (does not alter the
analysis specification) and a specification change (does alter it,
and requires re-locking) is principled but not always crisp in
practice. The ERSAF v1 PELT attribute-initialization fix
(\texttt{552cc0c}) is unambiguously an operational fix: a missing
attribute assignment that prevented the locked code from running
to completion. A hypothetical fix that changed a default parameter
value, however, would be harder to classify. The framework's
current resolution is to err on the side of re-locking whenever
the distinction is unclear, accepting the additional process cost
to preserve the artifact chain's interpretability. Future revisions
may formalize an intermediate category (``re-lock optional'')
governed by explicit criteria.

\paragraph{Scoring-rule contestability.}
AEPF's standing calibration requirement assumes that the underlying
probability estimator is amenable to standard scoring rules
\cite{brier1950, naeini2015ece, guo2017calibration}. Domains where
the appropriate scoring rule is contested---long-horizon
probabilistic forecasting under non-stationary processes,
open-ended generative evaluation, multi-modal output spaces---may
require methodology extensions beyond the framework's current
specification. The case studies in Section~\ref{sec:case-studies}
operate in domains where the scoring rules are well-established;
extension to contested-scoring domains is identified as work
beyond the present specification. Relatedly, the validation of the
calibration audit (Section~\ref{subsec:instrument-validation})
establishes soundness only against the parametric miscalibration
families it tested---temperature-scaled and systematic-logit-shift
perturbations---and does not establish that the instrument catches
every possible miscalibration. Difficulty-stratified miscalibration,
where a model is calibrated on easy items and broken on hard ones, is
untested; the soundness claim is bounded to the families validated
and is not a claim of completeness.

\paragraph{Reviewer-quality residual.}
Cross-model review under AEPF is only as strong as the reviewers
who perform it. A reviewer who fails to flag a methodological
error allows that error into the lock chain; a reviewer whose
review is itself shaped by undisclosed conflicts of interest
introduces bias the framework cannot detect. The framework's
mitigations---requiring reviewers to be institutionally separable
from the operator, requiring the review record to be committed
alongside the pre-registration, requiring reviewer Critical
findings to block lock---reduce but do not eliminate this residual.
The framework's claim is that an AEPF artifact chain is more
trustworthy than an equivalent un-locked methodology, not that
it is unconditionally trustworthy.

\section{Conclusion}
\label{sec:conclusion}

This paper has introduced the Avenridge Evaluation and Publication
Framework: a workflow discipline that integrates pre-registration,
cross-model review, materialization, lock-time commitment,
single-execution discipline, calibration assessment as a standing
requirement, constrained-interpretation reporting, and null-result
publication parity into a single operational framework anchored by
cryptographic commit references. The framework's novelty lies in the
integration and the artifact-level operationalization, not in any
individual component; pre-registration, blind analysis, calibration
assessment, and reproducibility frameworks each have substantial prior
literature, and AEPF inherits from each.

The four case studies presented in Section~\ref{sec:case-studies}
illustrate the framework's behavior across substantively different
evaluation contexts. The ERSAF cosmological audit demonstrates
AEPF's treatment of a NULL verdict under a locked full-lifecycle
study: a result that under conventional incentives might have
foregrounded a marginal Benjamini--Hochberg firing as a positive
finding instead reports the firing under two interpretive frames
with neither elevated above the other, with the methodology's
deliberate non-coverage of low-redshift peculiar-velocity
systematics documented in the forward-looking section. The
transformer calibration audits demonstrate the framework's
portability and the operational consequence of treating calibration
as a standing requirement: three audits, each under its own lock
chain, each surfacing a failure mode invisible to accuracy-only
evaluation. The Relational Coherence audit demonstrates the
framework on an open-weight language model under greedy decoding: a
high-framing-robustness regime named in advance as NULL-producing
returned that NULL through the full lock--calibrate--result
single-execution pipeline on commodity hardware. The Dynametrix-HRRR
independence diagnostic
demonstrates the viability-gate refinement
(Section~\ref{subsec:viability-gate}): a candidate study whose
informational yield was structurally over-determined was retired
before lock with a documented pre-specification disposition,
illustrating that AEPF's discipline extends backward to the
question of whether a study is worth locking at all. Beyond the case
studies, a methods-level validation against synthetic ground truth
establishes that the calibration audit used in the transformer
studies is itself sound---zero false-accept of broken calibration
across the tested miscalibration families---so that the NULL verdicts
above are evidence about the world rather than artifacts of the
instrument.

The framework's claim is procedural rather than substantive. AEPF
does not make any evaluation correct; it makes the relationship
between the question asked, the methodology applied, and the verdict
reported independently verifiable. The reader of an AEPF-locked
study verifies rather than infers; the operator who locks a
pre-registration commits to publish whatever the data return; the
cross-model reviewer leaves a written artifact that future readers
can examine. What changes is not what is true but what is
auditable.

Several directions for further work remain. Extension of the
calibration-assessment specification to domains where the
appropriate scoring rule is contested---long-horizon forecasting,
open-ended generative evaluation, multi-modal output spaces---is
identified in Section~\ref{sec:limitations} as work beyond the
present framework. Formalization of the operational-fix versus
specification-change boundary, including criteria for a possible
``re-lock optional'' intermediate category, would reduce ambiguity
in cases where the current discipline errs toward re-locking.
The viability-gate pattern introduced in
Section~\ref{subsec:viability-gate} is at present a refinement
under observation, with one operational application documented in
Section~\ref{subsec:dynametrix-hrrr-case}; further applications
across heterogeneous contexts will determine whether the trigger
conditions and constraints as currently stated hold up, or whether
they require revision before formalization as a full v1.0
component. Broader adoption of the framework would clarify the
empirical question of how AEPF discipline shifts the
published-result distribution---how the rate of null findings, the
rate of documented operational fixes, the rate of
constrained-interpretation reporting, and the rate of
pre-specification dispositions compare between AEPF-locked and
conventional evaluations. These remain open. What this paper has
presented is the framework itself, the artifacts that compose it,
and four case studies illustrating its application. If adopted
broadly, AEPF discipline would shift the rate of null-result
publication, the reproducibility of evaluation artifacts, and the
auditability of evaluation methodology in ways the framework
cannot itself measure; that measurement is itself an evaluation
question deserving AEPF discipline.

% =========================================================
% BIBLIOGRAPHY
% =========================================================

\bibliographystyle{unsrtnat}
\bibliography{references}

\end{document}
